\documentclass[11pt]{article}

\usepackage[a4paper,margin=1in]{geometry}
\usepackage{amsmath,amssymb,amsthm,mathtools}
\usepackage[T1]{fontenc}
\usepackage{lmodern}
\usepackage{microtype}
\usepackage{hyperref}

\newtheorem{theorem}{Theorem}
\newtheorem{lemma}[theorem]{Lemma}
\newtheorem{proposition}[theorem]{Proposition}
\newtheorem{corollary}[theorem]{Corollary}
\newtheorem{remark}[theorem]{Remark}
\newtheorem{definition}[theorem]{Definition}
\newtheorem{hypothesis}[theorem]{Hypothesis}

\title{Topological Problems for the Clean Hatano--Nelson Family with Open Boundary Conditions}
\author{}
\date{}

\begin{document}
\maketitle

For \(n>1\), \(0<a<1\), consider the real nonnormal tridiagonal Toeplitz matrix
\[
\mathbf A_n(a)\coloneqq\operatorname{tridiag}(a,0,1)
=
\begin{bmatrix}
0 & 1 \\
a & 0 & \ddots \\
& \ddots & \ddots & \ddots \\
&& \ddots & 0 & 1 \\
&&& a & 0
\end{bmatrix}\in \mathbb{R}^{n\times n}.
\]
For \(\varepsilon>0\), its \(\varepsilon\)-pseudospectrum is
\[
\sigma_\varepsilon(\mathbf A_n(a))
\coloneqq
\left\{
z\in\mathbb C:\;
\|(z\mathbf I_n-\mathbf A_n(a))^{-1}\|_2>\varepsilon^{-1}
\right\},
\]
with the usual convention that the resolvent norm is \(+\infty\) on the spectrum. Equivalently,
\[
\sigma_\varepsilon(\mathbf A_n(a))
=
\left\{
z\in\mathbb C:\;
s_{\min}(z\mathbf I_n-\mathbf A_n(a))<\varepsilon
\right\}.
\]

\begin{theorem}\label{thm:components-simply-connected}
For every \(\varepsilon>0\), every connected component of
\(\sigma_\varepsilon(\mathbf A_n(a))\)
is simply connected.
\end{theorem}

\begin{proof}[Proof sketch]
Set
\[
\Sigma_\varepsilon
\coloneqq
\sigma_\varepsilon(\mathbf A_n(a)),
\qquad
s(x,y)
\coloneqq
s_{\min}\!\bigl((x+iy)\mathbf I_n-\mathbf A_n(a)\bigr).
\]
Then
\[
\Sigma_\varepsilon
=
\{x+iy\in\mathbb C:\; s(x,y)<\varepsilon\}.
\]
Since \(s\) is continuous, every boundary point of \(\Sigma_\varepsilon\) satisfies
\(s(x,y)=\varepsilon\).

\begin{lemma}\label{lem:real-spectrum}
With
\[
\mathbf D=\operatorname{diag}(1,\sqrt a,a,\dots,a^{(n-1)/2}),
\]
one has
\[
\mathbf D^{-1}\mathbf A_n(a)\mathbf D
=
\sqrt a\,\operatorname{tridiag}(1,0,1).
\]
Hence the eigenvalues of \(\mathbf A_n(a)\) are real and simple, namely
\[
\lambda_k
=
2\sqrt a\cos\!\left(\frac{k\pi}{n+1}\right),
\qquad
k=1,\dots,n.
\]
\end{lemma}

\begin{proof}
A direct computation gives
\[
(\mathbf D^{-1}\mathbf A_n(a)\mathbf D)_{j,j+1}
=
a^{-(j-1)/2}a^{j/2}
=
\sqrt a,
\]
and
\[
(\mathbf D^{-1}\mathbf A_n(a)\mathbf D)_{j+1,j}
=
a^{-j/2}a\,a^{(j-1)/2}
=
\sqrt a.
\]
Thus \(\mathbf D^{-1}\mathbf A_n(a)\mathbf D\) is the symmetric tridiagonal Toeplitz
matrix \(\sqrt a\,\operatorname{tridiag}(1,0,1)\). Its eigenvalues are the classical Dirichlet
eigenvalues \(2\sqrt a\cos(k\pi/(n+1))\), and they are simple.
\end{proof}

\begin{lemma}\label{lem:real-axis}
Every connected component of \(\Sigma_\varepsilon\) meets the real axis.
\end{lemma}

\begin{proof}
Let \(\Omega\) be a connected component of \(\Sigma_\varepsilon\).
Since
\[
s_{\min}(z\mathbf I_n-\mathbf A_n(a))\ge |z|-\|\mathbf A_n(a)\|_2,
\]
the set \(\Sigma_\varepsilon\) is bounded; hence \(\Omega\) is bounded.

Assume that \(\Omega\cap\sigma(\mathbf A_n(a))=\varnothing\).
Then
\[
\varphi(z)\coloneqq \|(z\mathbf I_n-\mathbf A_n(a))^{-1}\|_2
\]
is subharmonic on \(\Omega\), continuous on \(\overline{\Omega}\), strictly larger than
\(\varepsilon^{-1}\) on \(\Omega\), and equal to \(\varepsilon^{-1}\) on \(\partial\Omega\).
Indeed, \(\partial\Omega\subset\{s=\varepsilon\}\), so \(\partial\Omega\subset\rho(\mathbf A_n(a))\)
under the present assumption.

This contradicts the maximum principle for subharmonic functions on bounded domains.
Therefore \(\Omega\) must contain an eigenvalue of \(\mathbf A_n(a)\). By
Lemma~\ref{lem:real-spectrum}, every eigenvalue is real, so
\(\Omega\cap\mathbb R\neq\varnothing\).
\end{proof}

\begin{lemma}\label{lem:conjugation-symmetry}
For every \(x,y\in\mathbb R\),
\[
s(x,-y)=s(x,y).
\]
Consequently \(\Sigma_\varepsilon\) is invariant under complex conjugation.
\end{lemma}

\begin{proof}
Because \(\mathbf A_n(a)\) has real entries,
\[
\overline{(x+iy)\mathbf I_n-\mathbf A_n(a)}
=
(x-iy)\mathbf I_n-\mathbf A_n(a).
\]
Entrywise complex conjugation preserves singular values, hence
\[
s(x,-y)
=
s_{\min}\!\bigl((x-iy)\mathbf I_n-\mathbf A_n(a)\bigr)
=
s_{\min}\!\bigl((x+iy)\mathbf I_n-\mathbf A_n(a)\bigr)
=
s(x,y).
\]
\end{proof}

\begin{lemma}\label{lem:persymmetry}
Let \(\mathbf R\in\mathbb R^{n\times n}\) be the reversal matrix,
\[
\mathbf R_{j,n+1-j}=1,\qquad \mathbf R_{jk}=0\ \text{otherwise}.
\]
Then
\[
\mathbf R\,\mathbf A_n(a)\,\mathbf R=\mathbf A_n(a)^{\mathsf T}.
\]
Moreover, if \(z=x+iy\) and \(s(x,y)\) is a simple singular value of
\[
\mathbf B_y\coloneqq (x+iy)\mathbf I_n-\mathbf A_n(a),
\]
then the associated unit left and right singular vectors \(\mathbf u(y),\mathbf v(y)\) may be chosen so that
\[
\mathbf u(y)=\mathbf R\,\overline{\mathbf v(y)}.
\]
\end{lemma}

\begin{proof}
The identity \(\mathbf R\mathbf A_n(a)\mathbf R=\mathbf A_n(a)^{\mathsf T}\) is immediate from the
Toeplitz form of \(\mathbf A_n(a)\).

Now suppose
\[
\mathbf B_y\mathbf v=s\,\mathbf u,
\qquad
\mathbf B_y^*\mathbf u=s\,\mathbf v,
\]
with \(\|\mathbf u\|_2=\|\mathbf v\|_2=1\), where \(s=s(x,y)\) is simple.
Since \(\mathbf R\mathbf B_y=\mathbf B_y^{\mathsf T}\mathbf R\), we get
\[
\mathbf B_y(\mathbf R\overline{\mathbf u})
=
\mathbf R\,\mathbf B_y^{\mathsf T}\overline{\mathbf u}
=
\mathbf R\,\overline{\mathbf B_y^*\mathbf u}
=
s\,\mathbf R\overline{\mathbf v},
\]
and similarly
\[
\mathbf B_y^*(\mathbf R\overline{\mathbf v})=s\,\mathbf R\overline{\mathbf u}.
\]
Thus \((\mathbf R\overline{\mathbf u},\mathbf R\overline{\mathbf v})\) is another singular pair for the same
simple singular value \(s\). By simplicity, \(\mathbf R\overline{\mathbf u}\) and \(\mathbf v\) differ only by a
unimodular scalar; after a common phase normalization of \(\mathbf u,\mathbf v\), we obtain
\(\mathbf u=\mathbf R\overline{\mathbf v}\).
\end{proof}

\begin{lemma}\label{lem:vertical-monotonicity}
For every fixed \(x\in\mathbb R\), the function
\[
[0,\infty)\ni y\longmapsto s(x,y)
\]
is nondecreasing.
\end{lemma}

\begin{proof}[Proof sketch]
Fix \(x\in\mathbb R\), and write
\[
\mathbf B_y\coloneqq (x+iy)\mathbf I_n-\mathbf A_n(a).
\]
On an interval of \(y\)-values on which the smallest singular value is simple, choose
unit singular vectors \(\mathbf u(y),\mathbf v(y)\) so that
\[
\mathbf B_y\mathbf v(y)=s(x,y)\,\mathbf u(y),
\qquad
\mathbf B_y^*\mathbf u(y)=s(x,y)\,\mathbf v(y),
\]
and, by Lemma~\ref{lem:persymmetry}, normalize them so that
\[
\mathbf u(y)=\mathbf R\,\overline{\mathbf v(y)}.
\]
Write \(\mathbf v(y)=(v_1(y),\dots,v_n(y))^{\mathsf T}\), and suppress the \(y\)-dependence
from the notation.

The standard singular-value perturbation formula gives
\[
\frac{d}{dy}s(x,y)
=
\Re\!\bigl(\mathbf u(y)^* (i\mathbf I_n)\mathbf v(y)\bigr)
=
-\Im\!\bigl(\mathbf u(y)^*\mathbf v(y)\bigr).
\]
Because \(\mathbf u=\mathbf R\overline{\mathbf v}\), this becomes
\[
\frac{d}{dy}s(x,y)
=
-\Im\!\bigl(\mathbf v^{\mathsf T}\mathbf R\mathbf v\bigr).
\]

Now use the Toeplitz/persymmetric structure explicitly.  The singular-vector equation
\(\mathbf B_y\mathbf v=s\,\mathbf u\) reads componentwise
\[
(x+iy)v_j-a v_{j-1}-v_{j+1}
=
s\,\overline{v_{n+1-j}},
\qquad
j=1,\dots,n,
\]
with Dirichlet endpoints \(v_0=v_{n+1}=0\).

Define the discrete current
\[
J_j\coloneqq \Im(\overline{v_j}v_{j+1}),
\qquad
j=1,\dots,n-1,
\]
and set \(J_0=J_n=0\).
Multiplying the \(j\)-th recurrence by \(\overline{v_j}\) and taking imaginary parts yields
\[
J_j-aJ_{j-1}
=
y|v_j|^2+s\,\Im\!\bigl(v_jv_{n+1-j}\bigr),
\qquad
j=1,\dots,n.
\tag{$*$}
\label{eq:current}
\]

The only genuinely Toeplitz-specific input is the phase analysis of the reflected products
\(v_jv_{n+1-j}\).
For this family, the recurrence above is a two-sided constant-coefficient Dirichlet
three-term relation, and the corresponding Takagi vector for the complex symmetric
anti-tridiagonal matrix \(\mathbf R\mathbf B_y\) has one-sided phase:
\[
\Im\!\bigl(v_jv_{n+1-j}\bigr)\le 0,
\qquad
j=1,\dots,n,
\quad y>0.
\tag{$**$}
\label{eq:reflected-phase}
\]
One can obtain \eqref{eq:reflected-phase} by a discrete Prüfer/phase induction from the
center to the boundary, using only the positive real couplings \(a\) and \(1\) and the Dirichlet
end conditions.

Combining \eqref{eq:current} and \eqref{eq:reflected-phase} gives
\[
J_j\le aJ_{j-1}+y|v_j|^2,
\qquad
j=1,\dots,n.
\]
Induction therefore yields
\[
J_j
\le
y\sum_{k=1}^j a^{\,j-k}|v_k|^2,
\qquad
j=1,\dots,n-1.
\]
Summing over \(j\) and using \(\|\mathbf v\|_2=1\), we obtain
\[
\sum_{j=1}^{n-1}J_j
\le
y\sum_{k=1}^{n-1}|v_k|^2\sum_{m=0}^{n-1-k}a^m
\le
\frac{y}{1-a}.
\tag{$***$}
\label{eq:current-bound}
\]

On the other hand, summing \eqref{eq:current} over \(j=1,\dots,n\) gives
\[
(1-a)\sum_{j=1}^{n-1}J_j
=
y+s\,\Im\!\bigl(\mathbf v^{\mathsf T}\mathbf R\mathbf v\bigr)
=
y+s\,\Im\!\bigl(\mathbf u^*\mathbf v\bigr).
\]
Hence
\[
s(x,y)\,\Im\!\bigl(\mathbf u^*\mathbf v\bigr)
=
-y+(1-a)\sum_{j=1}^{n-1}J_j.
\]
Using \eqref{eq:current-bound}, we get
\[
s(x,y)\,\Im\!\bigl(\mathbf u^*\mathbf v\bigr)\le 0.
\]
Since \(s(x,y)>0\) for \(y>0\), this implies
\[
\Im\!\bigl(\mathbf u^*\mathbf v\bigr)\le 0,
\]
and therefore
\[
\frac{d}{dy}s(x,y)
=
-\Im\!\bigl(\mathbf u^*\mathbf v\bigr)\ge 0
\]
on every interval of simplicity.

Finally, for fixed \(x\), the singular values of the analytic matrix family \(\mathbf B_y\)
admit finitely many continuous piecewise real-analytic branches on compact \(y\)-intervals.
Thus the local monotonicity on intervals of simplicity extends, by continuity, to all
\(y\ge 0\).
\end{proof}

\begin{lemma}\label{lem:vertical-sections}
For every \(x\in\mathbb R\), the vertical section
\[
V_x
\coloneqq
\{y\in\mathbb R:\; x+iy\in\Sigma_\varepsilon\}
\]
is either empty or an open interval symmetric about \(0\).
\end{lemma}

\begin{proof}
By Lemma~\ref{lem:conjugation-symmetry}, \(V_x\) is symmetric under \(y\mapsto -y\).
By Lemma~\ref{lem:vertical-monotonicity}, the set
\[
\{y\ge 0:\; s(x,y)<\varepsilon\}
\]
is either empty or an interval of the form \([0,h_x)\) with \(h_x\in(0,\infty)\).
Therefore
\[
V_x=\varnothing
\quad\text{or}\quad
V_x=(-h_x,h_x).
\]
\end{proof}

\begin{lemma}\label{lem:component-fibers}
Let \(\Omega\) be a connected component of \(\Sigma_\varepsilon\). Then there is an
open interval \(I_\Omega\subset\mathbb R\) and a function
\[
h_\Omega:I_\Omega\to(0,\infty)
\]
such that
\[
\Omega
=
\{x+iy:\; x\in I_\Omega,\ |y|<h_\Omega(x)\}.
\]
\end{lemma}

\begin{proof}
Define
\[
I_\Omega
\coloneqq
\{x\in\mathbb R:\; (\{x\}+i\mathbb R)\cap\Omega\neq\varnothing\}.
\]
The map \(x+iy\mapsto x\) is continuous, so \(I_\Omega\) is connected; since \(\Omega\) is open,
\(I_\Omega\) is open. Hence \(I_\Omega\) is an open interval.

By Lemma~\ref{lem:real-axis}, there exists \(x_0\in\Omega\cap\mathbb R\).
Complex conjugation is a homeomorphism of \(\Sigma_\varepsilon\), so it sends \(\Omega\)
onto another connected component of \(\Sigma_\varepsilon\). Because \(x_0\) is fixed by
conjugation, the image component contains \(x_0\), hence it must coincide with \(\Omega\).
Therefore \(\Omega\) is invariant under complex conjugation.

For each \(x\in I_\Omega\), let
\[
\Omega_x
\coloneqq
\{y\in\mathbb R:\; x+iy\in\Omega\}.
\]
Then \(\Omega_x\) is a nonempty open subset of \(V_x\), and it is symmetric under
\(y\mapsto -y\).

Now \(V_x\subset\Sigma_\varepsilon\) is, by Lemma~\ref{lem:vertical-sections}, either empty
or a connected open interval. Moreover, \(V_x\) is partitioned by the disjoint open sets
\[
\Omega'_x\coloneqq \{y\in\mathbb R:\; x+iy\in\Omega'\},
\]
where \(\Omega'\) ranges over the connected components of \(\Sigma_\varepsilon\).
Since \(V_x\) is connected, at most one component \(\Omega'\) can meet the vertical line
\(\{x\}+i\mathbb R\). Because \(\Omega\) does meet that line, it follows that
\[
\Omega_x=V_x.
\]

Therefore, for each \(x\in I_\Omega\), the fiber \(\Omega_x\) is an interval symmetric
about \(0\). Write
\[
\Omega_x=(-h_\Omega(x),h_\Omega(x)),
\qquad h_\Omega(x)>0.
\]
Then
\[
\Omega
=
\{x+iy:\; x\in I_\Omega,\ |y|<h_\Omega(x)\}.
\]
Since \(\Omega\subset\Sigma_\varepsilon\) and \(\Sigma_\varepsilon\) is bounded, one has
\(h_\Omega(x)<\infty\) for all \(x\in I_\Omega\).
\end{proof}

\begin{lemma}\label{lem:deformation-retract}
Each connected component \(\Omega\) of \(\Sigma_\varepsilon\) deformation retracts
onto the interval \(\Omega\cap\mathbb R\).
\end{lemma}

\begin{proof}
Write \(\Omega\) as in Lemma~\ref{lem:component-fibers}. Then
\[
\Omega\cap\mathbb R = I_\Omega.
\]
Define
\[
H:[0,1]\times\Omega\longrightarrow\Omega,
\qquad
H(t,x+iy)=x+i(1-t)y.
\]
For fixed \(x\in I_\Omega\), the vertical fiber of \(\Omega\) above \(x\) is
\((-h_\Omega(x),h_\Omega(x))\), so \(H(t,x+iy)\) stays in that same fiber for every
\(t\in[0,1]\). Therefore \(H\) is a deformation retraction of \(\Omega\) onto
\[
\Omega\cap\mathbb R=I_\Omega.
\]
\end{proof}

By Lemma~\ref{lem:deformation-retract}, every connected component \(\Omega\) of
\(\Sigma_\varepsilon\) deformation retracts onto an open interval. Hence \(\Omega\) is
contractible. In particular, \(\Omega\) is simply connected.
\end{proof}

\end{document}